\documentclass[11pt,a4paper]{article}

\usepackage[T1]{fontenc}
\usepackage[utf8]{inputenc}
\usepackage{lmodern}
\usepackage[margin=2.5cm]{geometry}
\usepackage{microtype}
\usepackage{xcolor}
\usepackage{graphicx}
\usepackage{array}
\usepackage{booktabs}
\usepackage{tabularx}
\usepackage{longtable}
\usepackage{enumitem}
\usepackage{hyperref}
\usepackage{titlesec}
\usepackage{fancyhdr}
\usepackage{setspace}
\usepackage{tikz}

\usetikzlibrary{positioning}

\definecolor{accent}{HTML}{25324A}
\definecolor{textgray}{HTML}{4D596A}
\definecolor{linegray}{HTML}{D9DEE7}
\hypersetup{
    colorlinks=true,
    linkcolor=accent,
    urlcolor=accent,
    citecolor=accent,
    pdftitle={Car Rental Company - Initial Project Report},
    pdfauthor={Project Team}
}

\setlength{\parindent}{0pt}
\setlength{\parskip}{0.55em}
\renewcommand{\arraystretch}{1.22}
\titleformat{\section}{\Large\bfseries\color{accent}}{}{0pt}{}
\titleformat{\subsection}{\normalsize\bfseries\color{accent}}{}{0pt}{}
\pagestyle{fancy}
\fancyhf{}
\fancyhead[R]{\footnotesize\color{textgray} Car Rental Company}
\fancyfoot[C]{\footnotesize\thepage}
\renewcommand{\headrulewidth}{0pt}

\renewcommand{\familydefault}{\sfdefault}

\newcommand{\projectTitle}{Car Rental Company}
\newcommand{\academicYear}{2025/2026}
\newcommand{\groupSection}{Informatics, Section 6}

\newcommand{\universityLogo}{polsl_logo.png}
\newcommand{\facultyLogo}{aei_logo.jpeg}

\newcommand{\logobox}[2]{%
  \begingroup
  \setlength{\fboxsep}{0pt}%
  \IfFileExists{#1}{\includegraphics[height=#2]{#1}}{%
    \fbox{\rule{0pt}{#2}\rule{0.2\linewidth}{0pt}\parbox[c][#2][c]{0.2\linewidth}{\centering Logo}}%
  }%
  \endgroup
}

\begin{document}
\begin{titlepage}
    \centering

    \begin{minipage}{0.45\linewidth}
        \centering
        \logobox{\universityLogo}{4.5cm}\\[0.4cm]
    \end{minipage}
    \hfill
    \begin{minipage}{0.45\linewidth}
        \centering
        \logobox{\facultyLogo}{3.5cm}\\[0.4cm]
    \end{minipage}

    \vspace{3.2cm}

    
\begin{tikzpicture}
        \node[draw=black!30, fill=black!5, rounded corners=10pt, inner sep=16pt, text width=0.9\linewidth, align=center] {
            {\LARGE \textbf{\projectTitle}} \\[0.4cm]
            {\Large Initial Project Report}
        };
    \end{tikzpicture}

    \vspace{3cm}

    \begin{flushleft}
    \setstretch{1.4}
    \textbf{Team members:} \\
    Tymoteusz Kłoska \\
    Szymon Molicki \\
    Marcjan Zawadzki \\
    Mateusz Znaleźniak \\
    Filipa Delfino \\[0.5cm]
    \textbf{Academic year:} \academicYear \\[0.2cm]
    \textbf{Section:} \groupSection \\[0.2cm]
    \textbf{Supervisor:} Hafedh Zghidi \\[0.5cm]
    \textbf{Date:} \today
    \end{flushleft}

    \vfill
    \centering

\end{titlepage}

\section{Project Overview}
The goal of the project is to design and implement a system for managing a car rental company. The application should support customer registration and profile editing, browsing available cars, making rental reservations, reviewing rental history, generating reports, online payments and using discount coupons. On the company side, the system should support fleet management, customer management, rental settlement, invoice generation and informing regular customers about promotions and discounts.

\section{Selected Tech Stack}
The project uses a Python-based stack. The selected technologies are presented in Table~\ref{tab:stack}.

\begin{table}[h!]
\centering
\small
\begin{tabularx}{\textwidth}{@{}p{2.7cm}p{4.2cm}X@{}}
\toprule
\textbf{Layer} & \textbf{Technology} & \textbf{Purpose} \\
\midrule
Frontend & React & User interface for customers and rental company staff. \\
Backend API & Python + FastAPI & REST API, business logic, validation and documentation. \\
ORM & SQLAlchemy & Object-relational mapping between Python models and the database. \\
Migrations & Alembic & Versioned database schema changes. \\
Database & PostgreSQL & Persistent relational data storage. \\
Containers & Docker Compose & Local development environment and service orchestration. \\
\bottomrule
\end{tabularx}
\caption{Selected project technology stack.}
\label{tab:stack}
\end{table}

\section{Initial Functional Scope}
\begin{table}[h!]
\centering
\small
\begin{tabularx}{\textwidth}{@{}p{2.8cm}X@{}}
\toprule
\textbf{Actor} & \textbf{Key functions} \\
\midrule
Customer & Register and edit personal data, browse available vehicles, create reservations, review rental history, pay online and use discounts. \\
Staff & Manage cars and customers, settle rentals, issue invoices and communicate promotions to regular customers. \\
\bottomrule
\end{tabularx}
\caption{Condensed functional scope.}
\end{table}

\section{Database Diagram}
The current database design is centered around the \texttt{RENTAL} table, which connects customers, cars, pickup and return locations and optional discounts.

\begin{figure}[h!]
    \centering
    \includegraphics[width=0.93\textwidth]{DB_diagram.pdf}
    \caption{Database diagram in table form.}
\end{figure}

\clearpage
\section{Core Entities}
\begin{table}[h!]
\centering
\small
\begin{tabularx}{\textwidth}{@{}p{2.7cm}X@{}}
\toprule
\textbf{Entity} & \textbf{Role in the system} \\
\midrule
\texttt{CUSTOMER} & Stores personal, contact, and driving license data required for rentals. \\
\texttt{LOCATION} & Represents rental branches used as car storage points and pickup or return places. \\
\texttt{CAR} & Stores vehicle identity, characteristics, status, rate, mileage and current location. \\
\texttt{DISCOUNT} & Stores discount codes and percentage-based promotions with validity dates. \\
\texttt{RENTAL} & Central transactional entity describing reservation and rental execution, cost calculation, invoice data and payment status. \\
\bottomrule
\end{tabularx}
\caption{Main entities visible in the current database diagram.}
\end{table}

\section{Conclusion}
This first report presents the agreed project direction in a concise form: the project scope, the selected Python-based technology stack and the initial relational database design.

\end{document}
